% !TeX root = main.tex
% Sample table, kept in its own file and included from chapter2.tex with
% % !TeX root = main.tex
% Sample table, kept in its own file and included from chapter2.tex with
% % !TeX root = main.tex
% Sample table, kept in its own file and included from chapter2.tex with
% % !TeX root = main.tex
% Sample table, kept in its own file and included from chapter2.tex with
% \input{table}. Separate files are optional but handy for large tables or
% tables generated by a script.
%
% GSO rules (the class handles the layout):
% - The title (\caption) goes ABOVE the table, in title case, with no final
%   period. It appears word for word in the List of Tables.
% - Tables are numbered by chapter (Table 2.1).
% - Table text is single-spaced; a small table must not split across pages.
% Refer to every table in the text with \ref{label}.

\begin{table}[htbp]
\centering
\caption{Major Differences Between Neural Networks and Cell Signaling Networks}
\label{table:NN_DIFF}

\begin{tabularx}{\linewidth}{|c|p{0.18\linewidth}|X|X|}
\hline
\textbf{Item} & \textbf{Salient Feature} & \textbf{Neural Networks} & \textbf{Cell Signaling Networks} \\
\hline
(i) & Nodes & All nodes typically identical & Nodes are not all equivalent in performance \\
\hline
(ii) & Structure & Typically layered with feed-forward connections only & Regulatory signals often give rise to feedback connections resulting in cycles \\
\hline
(iii) & Connectivity & Highly connected & More sparsely connected \\
\hline
(iv) & Learning Rule & Connectivity altered to perform a single function (e.g., via a back-propagation algorithm) & Cells must be able to respond to multiple stimuli effectively; changes typically occur via evolutionary processes \\
\hline
\end{tabularx}
\end{table}
. Separate files are optional but handy for large tables or
% tables generated by a script.
%
% GSO rules (the class handles the layout):
% - The title (\caption) goes ABOVE the table, in title case, with no final
%   period. It appears word for word in the List of Tables.
% - Tables are numbered by chapter (Table 2.1).
% - Table text is single-spaced; a small table must not split across pages.
% Refer to every table in the text with \ref{label}.

\begin{table}[htbp]
\centering
\caption{Major Differences Between Neural Networks and Cell Signaling Networks}
\label{table:NN_DIFF}

\begin{tabularx}{\linewidth}{|c|p{0.18\linewidth}|X|X|}
\hline
\textbf{Item} & \textbf{Salient Feature} & \textbf{Neural Networks} & \textbf{Cell Signaling Networks} \\
\hline
(i) & Nodes & All nodes typically identical & Nodes are not all equivalent in performance \\
\hline
(ii) & Structure & Typically layered with feed-forward connections only & Regulatory signals often give rise to feedback connections resulting in cycles \\
\hline
(iii) & Connectivity & Highly connected & More sparsely connected \\
\hline
(iv) & Learning Rule & Connectivity altered to perform a single function (e.g., via a back-propagation algorithm) & Cells must be able to respond to multiple stimuli effectively; changes typically occur via evolutionary processes \\
\hline
\end{tabularx}
\end{table}
. Separate files are optional but handy for large tables or
% tables generated by a script.
%
% GSO rules (the class handles the layout):
% - The title (\caption) goes ABOVE the table, in title case, with no final
%   period. It appears word for word in the List of Tables.
% - Tables are numbered by chapter (Table 2.1).
% - Table text is single-spaced; a small table must not split across pages.
% Refer to every table in the text with \ref{label}.

\begin{table}[htbp]
\centering
\caption{Major Differences Between Neural Networks and Cell Signaling Networks}
\label{table:NN_DIFF}

\begin{tabularx}{\linewidth}{|c|p{0.18\linewidth}|X|X|}
\hline
\textbf{Item} & \textbf{Salient Feature} & \textbf{Neural Networks} & \textbf{Cell Signaling Networks} \\
\hline
(i) & Nodes & All nodes typically identical & Nodes are not all equivalent in performance \\
\hline
(ii) & Structure & Typically layered with feed-forward connections only & Regulatory signals often give rise to feedback connections resulting in cycles \\
\hline
(iii) & Connectivity & Highly connected & More sparsely connected \\
\hline
(iv) & Learning Rule & Connectivity altered to perform a single function (e.g., via a back-propagation algorithm) & Cells must be able to respond to multiple stimuli effectively; changes typically occur via evolutionary processes \\
\hline
\end{tabularx}
\end{table}
. Separate files are optional but handy for large tables or
% tables generated by a script.
%
% GSO rules (the class handles the layout):
% - The title (\caption) goes ABOVE the table, in title case, with no final
%   period. It appears word for word in the List of Tables.
% - Tables are numbered by chapter (Table 2.1).
% - Table text is single-spaced; a small table must not split across pages.
% Refer to every table in the text with \ref{label}.

\begin{table}[htbp]
\centering
\caption{Major Differences Between Neural Networks and Cell Signaling Networks}
\label{table:NN_DIFF}

\begin{tabularx}{\linewidth}{|c|p{0.18\linewidth}|X|X|}
\hline
\textbf{Item} & \textbf{Salient Feature} & \textbf{Neural Networks} & \textbf{Cell Signaling Networks} \\
\hline
(i) & Nodes & All nodes typically identical & Nodes are not all equivalent in performance \\
\hline
(ii) & Structure & Typically layered with feed-forward connections only & Regulatory signals often give rise to feedback connections resulting in cycles \\
\hline
(iii) & Connectivity & Highly connected & More sparsely connected \\
\hline
(iv) & Learning Rule & Connectivity altered to perform a single function (e.g., via a back-propagation algorithm) & Cells must be able to respond to multiple stimuli effectively; changes typically occur via evolutionary processes \\
\hline
\end{tabularx}
\end{table}
